\chapter{Kết luận và hướng phát triển}

\section{Kết quả đạt được}

Khóa luận này đã trình bày một quy trình xây dựng đồ thị tri thức từ dữ liệu tiếng Việt về Trường Đại học Công nghệ. Quy trình bắt đầu từ dữ liệu công khai được thu thập từ các trang web của nhà trường, sau đó xử lý thành dữ liệu văn bản có cấu trúc ổn định hơn, trích xuất các đối tượng và quan hệ quan trọng, xử lý các cách gọi đồng tham chiếu của cùng một đối tượng, cập nhật lại quan hệ và giữ lại bằng chứng nguồn để phục vụ kiểm chứng. Trên nền đồ thị đã xây dựng, khóa luận triển khai các kiến trúc truy vấn phục vụ ứng dụng hỏi đáp nhằm kiểm tra giá trị sử dụng của nguồn tri thức có cấu trúc.

Về mặt xây dựng đồ thị tri thức, khóa luận đã đề xuất và triển khai một quy trình từ đầu đến cuối cho dữ liệu văn bản tiếng Việt phi cấu trúc. Kết quả của quy trình không chỉ là một tập các thực thể và quan hệ rời rạc, mà là một nguồn tri thức có tổ chức, có định danh thực thể ổn định hơn và có thông tin truy vết nguồn gốc. Việc giữ lại nguồn thông tin giúp các quan hệ trong đồ thị có thể được đối chiếu lại khi cần, đồng thời tạo điều kiện cho bước đánh giá chất lượng ở lớp dữ liệu tri thức.

Về mặt chất lượng đồ thị, khóa luận đã trình bày quy trình xử lý đồng tham chiếu thực thể nhằm giảm tình trạng phân mảnh do cùng một đối tượng xuất hiện dưới nhiều cách gọi khác nhau. Quy trình này giúp gom các thực thể có cùng tham chiếu, tạo thực thể chuẩn và cập nhật lại các quan hệ liên quan. Kết quả đánh giá trên 500 quan hệ được chọn ngẫu nhiên cho thấy 93,4\% quan hệ được bằng chứng trong dữ liệu nguồn hỗ trợ. Kết quả này cho thấy đồ thị tri thức được xây dựng có mức độ đáng tin cậy tương đối tốt trong phạm vi dữ liệu được khảo sát.

Về mặt ứng dụng hỏi đáp, khóa luận đã sử dụng đồ thị tri thức như một nguồn bằng chứng có cấu trúc bên cạnh văn bản gốc. Các kiến trúc truy vấn văn bản, truy vấn đồ thị trực tiếp, truy vấn đồ thị sâu và truy vấn kết hợp cho phép so sánh nhiều cách sử dụng bằng chứng khác nhau. Kết quả ở Chương~6 cho thấy kiến trúc truy vấn đồ thị sâu đạt điểm bám bằng chứng và điểm đúng câu trả lời cao nhất, trong khi kiến trúc truy vấn kết hợp đạt điểm liên quan câu trả lời cao nhất. So với kiến trúc truy vấn văn bản, truy vấn đồ thị sâu cải thiện khoảng 11,5\% ở điểm bám bằng chứng và 15,9\% ở điểm đúng câu trả lời; truy vấn kết hợp cải thiện khoảng 11,7\% ở điểm liên quan câu trả lời.

Nhìn chung, khóa luận đạt được ba kết quả chính. Thứ nhất, khóa luận đề xuất một quy trình hoàn chỉnh để biến nguồn văn bản phân tán thành đồ thị tri thức có thể kiểm chứng. Thứ hai, khóa luận xây dựng một bộ dữ liệu đồ thị tri thức phục vụ hỏi đáp trong phạm vi Trường Đại học Công nghệ. Thứ ba, khóa luận tổ chức đánh giá theo hai lớp, gồm đánh giá chất lượng quan hệ trong đồ thị và đánh giá tác động của đồ thị ở lớp ứng dụng hỏi đáp.

\section{Hạn chế}

Mặc dù đã hoàn thiện các thành phần chính của quy trình, khóa luận vẫn còn một số hạn chế.

Thứ nhất, quy mô dữ liệu thô còn hạn chế so với nhu cầu bao phủ đầy đủ thông tin của một trường đại học. Tập dữ liệu hiện tại chủ yếu dựa trên các nguồn công khai đã thu thập được, do đó có thể chưa phản ánh hết các thay đổi mới, các tài liệu lịch sử, các văn bản nội bộ hoặc các nội dung có cấu trúc đặc thù. Khi dữ liệu đầu vào chưa đủ rộng, đồ thị tri thức cũng có thể thiếu một số thực thể, quan hệ hoặc ngữ cảnh cần thiết cho hỏi đáp.

Thứ hai, khóa luận chưa xây dựng được tập dữ liệu chuẩn có nhãn thủ công đủ lớn để đánh giá đầy đủ độ chính xác của bước trích xuất tri thức và bước xử lý đồng tham chiếu thực thể. Vì vậy, các giai đoạn này mới được báo cáo thông qua thống kê đầu ra, ví dụ minh họa và đánh giá độ đúng theo bằng chứng trên mẫu quan hệ, thay vì đánh giá toàn diện trên toàn bộ đồ thị.

Thứ ba, mô hình nhúng được sử dụng là mô hình đa ngôn ngữ tổng quát, chưa được tối ưu riêng cho miền dữ liệu của Trường Đại học Công nghệ. Điều này có thể ảnh hưởng đến các bước dựa trên độ tương đồng ngữ nghĩa như truy xuất văn bản, phân cụm ứng viên đồng tham chiếu, liên kết thực thể và tìm kiếm ngữ nghĩa trong ứng dụng hỏi đáp.

Thứ tư, chất lượng của nhiều thành phần phụ thuộc vào sức mạnh và độ ổn định của mô hình ngôn ngữ lớn. Các bước trích xuất tri thức, hợp nhất thực thể, sinh câu trả lời và đánh giá tự động đều có thể thay đổi theo mô hình được chọn, cách nhắc lệnh, giới hạn ngữ cảnh và chính sách suy luận của mô hình. Vì vậy, việc lựa chọn mô hình cần được thực hiện cẩn thận, đặc biệt khi hệ thống được mở rộng sang môi trường vận hành thực tế.

Thứ năm, khóa luận đi theo hướng kiểm soát kiểu thực thể và quan hệ để giữ đồ thị tri thức ổn định, nhưng cách tiếp cận này có thể bỏ sót một số kiểu thực thể phát sinh ngoài lược đồ ban đầu. Đối với miền dữ liệu đại học, một bản thể học tốt cần có sự tham gia của chuyên gia miền để phân tích các nhóm thực thể, quan hệ và ràng buộc đủ bao phủ nhưng vẫn tránh làm lược đồ trở nên quá phức tạp.

Thứ sáu, phần đánh giá hỏi đáp được thực hiện trên 146 câu hỏi và sử dụng bộ chấm tự động. Cách đánh giá này giúp so sánh các kiến trúc truy vấn trong cùng một giao thức, nhưng chưa thay thế cho đánh giá thủ công độc lập bởi con người. Ngoài ra, khóa luận chưa thực hiện kiểm định thống kê theo từng câu hỏi để xác nhận mức ý nghĩa của chênh lệch giữa các kiến trúc truy vấn.

Thứ bảy, tốc độ trả lời của các kiến trúc truy vấn đồ thị sâu và truy vấn kết hợp còn chậm. Nguyên nhân có thể đến từ kiến trúc truy vấn nhiều bước, số lượt gọi mô hình ngôn ngữ lớn, thao tác truy xuất và phần mã triển khai chưa được tối ưu sâu. Đây chưa phải trọng tâm chính của khóa luận, nhưng là hạn chế quan trọng nếu muốn triển khai hệ thống hỏi đáp cho người dùng cuối.

\section{Hướng phát triển}

Từ các hạn chế trên, khóa luận có thể được phát triển tiếp theo một số hướng.

Thứ nhất, cần mở rộng quy mô và tính cập nhật của dữ liệu. Một hướng phát triển thực tế là xây dựng hệ thống hoạt động gần thời gian thực cho Trường Đại học Công nghệ, trong đó dữ liệu mới được tự động thu thập, phát hiện thay đổi, xử lý, trích xuất tri thức, cập nhật đồ thị và ghi nhận nguồn gốc. Cách tổ chức này giúp đồ thị tri thức phản ánh kịp thời các thay đổi về đơn vị, nhân sự, chương trình đào tạo, thông báo và sự kiện.

Thứ hai, cần xây dựng bộ dữ liệu đánh giá có nhãn cho các bước trích xuất tri thức và xử lý đồng tham chiếu thực thể. Bộ dữ liệu này nên bao gồm nhãn về thực thể, quan hệ và các cặp thực thể cùng tham chiếu một đối tượng, từ đó cho phép đánh giá chính xác hơn chất lượng của từng giai đoạn trong quy trình xây dựng đồ thị.

Thứ ba, có thể tinh chỉnh hoặc huấn luyện bổ sung mô hình nhúng riêng cho miền dữ liệu của Trường Đại học Công nghệ. Mô hình nhúng chuyên biệt có thể cải thiện truy xuất ngữ nghĩa, phát hiện thực thể đồng tham chiếu, liên kết câu hỏi với thực thể trong đồ thị và lựa chọn ngữ cảnh cho ứng dụng hỏi đáp.

Thứ tư, cần tiếp tục hoàn thiện bản thể học cho miền dữ liệu đại học. Việc này nên có sự tham gia của chuyên gia miền để xác định các kiểu thực thể, kiểu quan hệ, thuộc tính và ràng buộc quan trọng. Một bản thể học đủ bao phủ sẽ giúp giảm bỏ sót thực thể phát sinh, đồng thời làm cho kết quả trích xuất và hợp nhất thực thể nhất quán hơn.

Thứ năm, cần hoàn thiện cơ chế human fallback và kiểm duyệt thủ công có hỗ trợ. Hệ thống có thể gợi ý nhiều ứng viên thực thể trùng lặp hoặc đồng tham chiếu hơn, cho phép người dùng hợp nhất, tách, chỉnh sửa thuộc tính, sửa quan hệ và ghi nhận lý do hiệu chỉnh. Các quyết định này có thể được lưu lại như dữ liệu phản hồi để cải thiện quy trình xử lý đồng tham chiếu trong các lần chạy sau.

Thứ sáu, cần nâng cấp kiến trúc hỏi đáp theo hướng tối ưu tốc độ. Các hướng có thể xem xét gồm giảm số lượt gọi mô hình ngôn ngữ lớn, cache kết quả truy xuất, tối ưu liên kết thực thể, thu hẹp phạm vi truy vấn đồ thị, xử lý song song các bước độc lập và thiết kế lại pipeline deep graph search hoặc hybrid search để phù hợp hơn với yêu cầu tương tác thực tế.

Thứ bảy, cần mở rộng đánh giá lớp hỏi đáp bằng đánh giá thủ công độc lập, phân tích lỗi theo từng nhóm câu hỏi và bổ sung các chỉ số phía truy xuất. Việc này giúp làm rõ câu hỏi nào phù hợp với semantic search, câu hỏi nào cần deep graph search, và câu hỏi nào nên dùng hybrid search giữa văn bản với đồ thị.

\section{Kết luận chung}

Khóa luận đã xây dựng một quy trình tạo đồ thị tri thức từ dữ liệu văn bản tiếng Việt và kiểm chứng khả năng sử dụng đồ thị trong ứng dụng hỏi đáp. Điểm trọng tâm của khóa luận không nằm ở việc đề xuất một phương pháp hỏi đáp hoàn toàn mới, mà nằm ở cách xây dựng, làm sạch, truy vết và đánh giá nguồn tri thức trước khi đưa vào khai thác.

Kết quả thực nghiệm cho thấy đồ thị tri thức có giá trị như một nguồn bằng chứng có cấu trúc. Khi được truy vấn đúng cách, đồ thị giúp câu trả lời bám bằng chứng hơn và đúng với đáp án tham chiếu hơn so với kiến trúc chỉ dựa trên văn bản. Tuy nhiên, lợi ích này không tự động xuất hiện ở mọi cách sử dụng đồ thị; nó phụ thuộc vào chất lượng trích xuất, xử lý đồng tham chiếu, liên kết thực thể và chiến lược truy vấn. Vì vậy, các hướng phát triển tiếp theo cần tiếp tục ưu tiên chất lượng đồ thị, khả năng kiểm chứng dữ liệu và hiệu quả của các kiến trúc truy vấn ở lớp ứng dụng hỏi đáp.
